\section{质数筛}

\verb|Sieve()| 预处理 $[0,N)$ 内每个数是否为质数，\texttt{np[i]} 为 \texttt{true} 表示 $i$ 是合数；$N$ 在编译期确定，一般取 $10^7$ 量级。本机运行 $50$ 次的平均耗时（毫秒）：

\begin{center}
\begin{tabular}{lcc}
\hline
数据规模 & 线性筛(ms) & 埃氏筛(ms) \\
\hline
$2\times 10^6$ & $8.354$ & $8.812$ \\
$1\times 10^7$ & $41.562$ & $44.417$ \\
$2\times 10^7$ & $83.388$ & $115.250$ \\
$5\times 10^7$ & $205.396$ & $308.354$ \\
$1\times 10^8$ & $422.438$ & $673.562$ \\
\hline
\end{tabular}
\end{center}

埃氏筛实现简单，适合 $N$ 较小（约 $10^7$）的场合：

\begin{lstlisting}
constexpr int N = 1e7 + 3;
std::bitset<N> np;
void Sieve(){
    np[0] = np[1] = 1;
    for (int i = 2;i < N;i++){
        if (np[i]) continue;
        for (int j = i << 1;j < N;j += i){
            np[j] = 1;
        }
    }
}
\end{lstlisting}

欧拉筛（线性筛）常数更优，并顺带得到质数表 \texttt{primes}：

\begin{lstlisting}
constexpr int N = 1e7 + 3;
std::bitset<N> np;
std::vector<int> primes;
void Sieve(){
    np[0] = np[1] = 1;
    for (int i = 2;i < N;i++){
        if (!np[i]) primes.push_back(i);
        for (const int& x : primes){
            if (x * i >= N) break;
            np[x * i] = 1;
            if (i % x == 0) break;
        }
    }
}
\end{lstlisting}

\section{Miller--Rabin 素性测试}

\verb|isPrime(n)| 判断 $64$ 位整数 $n$ 是否为质数，固定基底、确定性结果。辅助函数：

\begin{itemize}
    \item \verb|quickPower(a,p,m)|：取模快速幂。
    \item \verb|mul(x,y,m)|：借助 \texttt{\_\_int128} 的不溢出取模乘法。
\end{itemize}

\begin{lstlisting}
using i64 = long long;
i64 mul(const i64& x,const i64& y,const i64& m){
    return static_cast<__int128>(x) * y % m;
}

i64 quickPower(i64 a,i64 p,i64 m){
    i64 r = 1;
    for(a %= m;p;p>>=1,a = mul(a,a,m)) {
        if (p & 1) r = mul(r,a,m);
    }
    return r;
}

bool isPrime(i64 n){
    if (n == 2 || n == 3) return 1;
    if (n % 2 == 0 || n % 3 == 0 || n < 2) return 0;
    int p2 = __builtin_ctzll(n-1);
    constexpr std::array<int,7> P = {2,325,9375,28178,450775,9780504,1795265022};
    for (const auto& p : P){
        if (p % n == 0) break;
        i64 x = quickPower(p, (n-1) >> p2, n);
        if (x == 1 || x == n-1) continue;
        for (int i = 1;i <= p2;i++){
            i64 y(x);
            x = mul(x,x,n);
            if (x == 1) {
                if (y != n-1) return 0;
                break;
            }
        }
        if (x != 1) return 0;
    }
    return 1;
}
\end{lstlisting}

\section{质因数分解（Pollard--Rho）}

前置模板：\textbf{\texttt{Miller--Rabin 素性测试}}。

\verb|factorize(v)| 返回 $v$ 的全部质因子（可重复、升序排列），要求 $v\ge 1$，期望复杂度 $\mathcal{O}(v^{\frac14})$。

\begin{lstlisting}
using i64 = long long;

std::vector<i64> factorize(i64 v) {
    static constexpr int C = 1e4;
    std::vector<i64> p;
    while ((v & 1) == 0) {
        p.push_back(2);
        v >>= 1;
    }
    std::function<void(i64)> PollardRho = [&](i64 n)->void {
        if (isPrime(n)) {
            p.push_back(n);
            return;
        }
        if (n <= C) {
            for (int i = 2;i * i <= n;i++){
                while (n % i == 0) {
                    p.push_back(i);
                    n /= i;
                }
            }
            if (n > 1) p.push_back(n);
            return;
        }
        static std::mt19937_64 rng(time(0));
        std::uniform_int_distribution<i64> uf(1,n-1);
        i64 c = uf(rng);
        auto f = [&](i64 x)->i64 {
            return (static_cast<__int128>(x) * x + c) % n;
        };
        while (1) {
            i64 x = 0,y = 0;
            i64 d = 1,prod = 1;
            int lmt = 1,cnt = 0;
            while (d == 1) {
                y = f(y);
                prod = mul(prod,std::abs(x-y),n);
                ++cnt;
                if (cnt % 127 == 0) {
                    d = std::gcd(prod,n);
                    prod = 1;
                }
                if (lmt == cnt) {
                    x = y;
                    cnt = 0;
                    lmt <<= 1;
                    d = std::gcd(prod,n);
                    prod = 1;
                }
            }
            if (d != n) {
                PollardRho(n/d);
                PollardRho(d);
                return;
            }
            c = uf(rng);
        }
    };
    PollardRho(v);
    std::sort(p.begin(),p.end());
    return p;
}
\end{lstlisting}
